% Genereller Header
\documentclass[10pt,twoside,a4paper,fleqn]{article}
\usepackage[left=1cm,right=1cm,top=1cm,bottom=1cm,includeheadfoot]{geometry}
\usepackage[english]{babel}
\usepackage[utf8]{inputenc}

% Pakete
\usepackage{amssymb}
\usepackage{amsmath}
\usepackage{bm}       % Bold Math
\usepackage{fancybox}
\usepackage{graphicx}
\usepackage{color}
\usepackage{lastpage}
\usepackage{wrapfig}
\usepackage{fancyhdr}
\usepackage{hyperref}
\usepackage{verbatim}
\usepackage{pdflscape} % landscape
\usepackage{multirow} % zellen in tabellen verbinden
\usepackage{multicol} 
%\usepackage{slashbox} % getrennte zelle in tabelle
\usepackage{enumitem}
% \usepackage{array} % anordnung in tabellen
\usepackage{tikz}
\usepackage{pgfplots}
\usetikzlibrary{pgfplots.groupplots}


%%%%%%%%%%%%%%%%%%%%
% Generelle Makros %
%%%%%%%%%%%%%%%%%%%%
\newcommand{\formelbuch}[1]{$_{\textcolor{red}{\mbox{\small{p. #1}}}}$}
\newcommand{\verweis}[2]{ {\small (siehe auch \ref{#1}, #2 (S. \pageref{#1}))
}}
\newcommand{\subsubadd}[1]{\textcolor{black}{\mbox{#1}}}
\newenvironment{liste}[0]{
	\begin{list}{$\bullet$}{\setlength{\itemsep}{0cm}\setlength{\parsep}{0cm} \setlength{\topsep}{0cm}}}
    {\end{list}}

\newenvironment{beschreibListe}[0]{\vspace{-2mm}
	\begin{description}[noitemsep, nolistsep]\setlength{\itemsep}{0cm} \setlength{\parsep}{0cm} \setlength{\topsep}{0cm} \setlength{\parskip}{0cm}}
    {\end{description}}

    
\newcommand{\logd}[0]{\log_{10}}
\newcommand{\subsubsubsection}[1]{\textbf{#1}}

\newenvironment{aufzaehlung}[0]{
	\begin{enumerate}\setlength{\itemsep}{0cm}\setlength{\parsep}{0cm} \setlength{\topsep}{0cm} \setlength{\parskip}{0cm}} {\end{enumerate}}

\newcommand{\abbHeight}[3]{
	\begin{center}
		\includegraphics[height=#2]{./bilder/#1} \\
		#3
    \end{center}
}

\newcommand{\hilight}[1]{\colorbox{yellow}{#1}}
\newcommand{\todo}[1]{\colorbox{red}{TODO: #1}}

\newcommand{\skriptsection}[2]{\section{\texorpdfstring{#1 \formelbuch{#2}}{#1}}}
\newcommand{\skriptsubsection}[2]{\subsection{\texorpdfstring{#1 \formelbuch{#2}}{#1}}}
\newcommand{\skriptsubsubsection}[2]{\subsubsection{\texorpdfstring{#1 \formelbuch{#2}}{#1}}}
\newcommand{\skriptsubsubsubsection}[2]{\subsubsubsection{\texorpdfstring{#1 \formelbuch{#2}}{#1}}}

%%%%%%%%%%
% Farben %
%%%%%%%%%%
\definecolor{black}{rgb}{0,0,0}
\definecolor{red}{rgb}{1,0,0}
\definecolor{white}{rgb}{1,1,1}
\definecolor{grey}{rgb}{0.8,0.8,0.8}

%%%%%%%%%%%%%%%%%%%%%%%%%%%%
% Mathematische Operatoren %
%%%%%%%%%%%%%%%%%%%%%%%%%%%%
\DeclareMathOperator{\sinc}{sinc}
\DeclareMathOperator{\sgn}{sgn}
\newcommand{\error}{\mathrm{error}}
\newcommand{\bfx}{\mathbf{x}}
\DeclareMathOperator{\expected}{\mathcal{E}}


% Fouriertransformationen
\unitlength1cm
\newcommand{\FT}
{
\begin{picture}(1,0.5)
\put(0.2,0.1){\circle{0.14}}\put(0.27,0.1){\line(1,0){0.5}}\put(0.77,0.1){\circle*{0.14}}
\end{picture}
}


\newcommand{\IFT}
{
\begin{picture}(1,0.5)
\put(0.2,0.1){\circle*{0.14}}\put(0.27,0.1){\line(1,0){0.45}}\put(0.77,0.1){\circle{0.14}}
\end{picture}
}



%%%%%%%%%%%%%%%%%%%%%%%%%%%%
% Allgemeine Einstellungen %
%%%%%%%%%%%%%%%%%%%%%%%%%%%%
%pdf info
\hypersetup{pdfauthor={\authorinfo},pdftitle={\titleinfo},colorlinks=false}
\author{\authorinfo}
\title{\titleinfo}

%Kopf- und Fusszeile
\pagestyle{fancy}
\fancyhf{}
%Linien oben und unten
\renewcommand{\headrulewidth}{0.5pt} 
\renewcommand{\footrulewidth}{0.5pt} 

\fancyhead[L]{\titleinfo{ }- Summary}
%Kopfzeile rechts bzw. aussen
\fancyhead[R]{\today{ }- Page \thepage/\pageref{LastPage}}
\fancyfoot[L]{\includegraphics{images/ccbyncsa.png}}
\fancyfoot[R]{\authorinfo}

% Einr�cken verhindern versuchen
\setlength{\parindent}{0pt}

\setcounter{tocdepth}{2}
